\documentclass{article}
\usepackage[utf8]{inputenc}
\usepackage{amsmath}
\usepackage{amsfonts}
\usepackage{amssymb}
\usepackage{siunitx}
\usepackage[left=1.0cm, right=1.0cm, top=1cm, bottom=1cm]{geometry}
\begin{document}
	\twocolumn
	\section{Das SI-System}
	
	\begin{tabular}{|l|l|l|l|}
		\hline\tiny
		\textbf{Zehnerpotenz} & \textbf{Vorsilbe} & \textbf{Kurzzeichen} & \textbf{Beispiel} \\
		\hline
		$10^{18}$ & Exa & E & EV \\
		$10^{15}$ & Peta & P & PV \\
		$10^{12}$ & Tera & T & TV \\
		$10^{9}$ & Giga & G & GHz \\
		$10^{6}$ & Mega & M & MHz \\
		$10^{3}$ & Kilo & k & K$\Omega$ \\
		$10^{2}$ & Hekto & h & hl \\
		$10^{1}$ & Deka & da & daA \\
		$10^{0}$ & & & \\
		$10^{-1}$ & Dezi & d & dm \\
		$10^{-2}$ & Zenti & C & cm \\
		$10^{-3}$ & Milli & m & mA \\
		$10^{-6}$ & Mikro & $\mu$ & $\mu$A \\
		$10^{-9}$ & Nano & n & nV \\
		$10^{-12}$ & Piko & p & pF \\
		$10^{-15}$ & Femto & f & fF \\
		$10^{-18}$ & Atto & a & aF \\
		\hline
	\end{tabular}
	
	\section{Elektrisches Feld \& Kapazität}
	
	Bei mehreren Feldern:
	Feldstärke $\vec{E}$ am Ort für die Ladung $Q_{0}$ im Ort $\vec{r_{0}}:$
	$\vec{E}(\vec{r})=\frac{Q_{0}}{4\pi\epsilon\cdot||\vec{r}-\vec{r_{0}}||}\cdot\frac{\vec{r}-\vec{r_{0}}}{||\vec{r}-\vec{r_{0}}||}$
	$E_{res}(\vec{r})=\sum_{i}E_{i}(\vec{r})$
	
	Homogene elektrische Feldstärke E
	$E=\frac{U}{d}=\frac{Q}{\epsilon_{0}\epsilon_{r}A}$ $\left[\SI{1}{\volt\per\meter}\right]$
	mit $d=$ Plattenabstand, $\epsilon_{0}=8.854187817 \times 10^{-12}$
	\SI{}{\volt\per\meter}
	
	\section{Elektrische Energie \& Leistung}
	Elektrische Leistung P:
	$P=U\cdot I=\frac{U^{2}}{R}=I^{2}\cdot R$ $\left[\SI{1}{\watt}=\SI{1}{\joule\per\second}=\SI{1}{\volt\ampere}=\SI{1}{\kilogram\meter\squared\per\second\cubed}\right]$
	Elektrische Energie $E_{el}$:
	$E_{el}=U\cdot I\cdot t$ $\left[\SI{1}{\watt\second}=\SI{1}{\joule}=\SI{1}{\newton\meter}=\SI{1}{\volt\ampere\second}=\SI{1}{\volt}\right]$
	Elektrische Kapazität
	$C=\frac{Q}{U}$
	$\left[\SI{1}{\farad}=\SI{1}{\ampere\second\per\volt}\right]$
	Parallelschaltung von Kondensatoren:
	$C_{ges}=C_{1}+C_{2}+\cdot\cdot\cdot+C_{n}$
	Reihenschaltung von Kondensatoren:
	$\frac{1}{C_{ges}}=\frac{1}{C_{1}}+\frac{1}{C_{2}}+\cdot\cdot\cdot+\frac{1}{C_{n}}$
	Kinetische Energie $E_{kin}$
	$E_{kin}=\frac{1}{2}\cdot m\cdot v^{2}$
	$\left[\SI{1}{\joule}\right]$
	Potentielle Energie $E_{pot}$
	$E_{pot}=m\cdot g\cdot h$
	$\left[\SI{1}{\joule}\right]$
	Thermische Energie $E_{th}$
	$E_{th}=m\cdot c\cdot\Delta T$
	$\left[\SI{1}{\joule}\right]$
	mit m Masse, $c=$ spezifische Wärmekapazität und $\triangle T=$ Temperaturdifferenz
	Tabelle 1: Umrechnung von Energieeinheiten
	\begin{tabular}{|l|l|l|l|}
		\hline
		& \multicolumn{3}{c|}{\textbf{Umrechnung in}} \\
		\textbf{} & \textbf{J, Ws} & \textbf{kcal} & \textbf{kWh} \\
		\hline
		\textbf{J, Ws} & 1 & $2.39\cdot10^{-4}$ & $2.78\cdot10^{-7}$ \\
		\textbf{kcal} & 4186 & 1 & $1.16\cdot10^{-3}$ \\
		\textbf{kWh} & $3.6\cdot10^{6}$ & 860 & 1 \\
		\hline
	\end{tabular}
	
	\section{Elektrischer Widerstand}
	
	Elektrischer Widerstand $R$:
	$R=\frac{U}{I}$
	$\left[\SI{1}{\ohm}\right]$
		Elektrische Widerstände Parallel:
	$$R_{ges}=\frac{1}{\frac{1}{R_{1}}+\frac{1}{R_{2}}+...+\frac{1}{R_{n}}}$$
	\section{Auf- und Entladung eines Kondensators}
	
	Feldenergie im Kondensator:
	$E=\frac{1}{2}\cdot C\cdot U^{2}$
	
	Aufladung:
	$U_{C}(t)=U_{0}\cdot(1-e^{-\frac{t}{\tau}})=U_{0}\cdot(1-e^{-\frac{t}{R_{C}\cdot C}})$
	$I_{C}(t)=\frac{U_{0}}{R_{C}}\cdot e^{-\frac{t}{\tau}}=I_{0}\cdot e^{-\frac{t}{R_{C}\cdot C}}$
	
	Während des Entladevorgangs eines Kondensators gilt:
	$U_{C}(t)=U_{0}\cdot e^{-\frac{t}{\tau}}=U_{0}\cdot e^{-\frac{t}{R_{C}\cdot C}}$
	$I_{C}(t)=-\frac{U_{0}}{R_{C}}\cdot e^{-\frac{t}{\tau}}=-I_{0}\cdot e^{-\frac{t}{R_{C}\cdot C}}$
	
	Um die Lade-/Entladezeit bis zu einer bestimmten Spannung zu bestimmen, stellt man die Lade- / Entladeformeln nach $t$ um.
	$-\tau\cdot \ln\left(1-\frac{U}{U_{0}}\right)=t$
	
	Elektrische Widerstände in Reihe:
	$R_{ges}=R_{1}+R_{2}+...+R_{n}$
	
	mit
	$\tau=R_{C}\cdot C$
	
	\section{Magnetisches Feld \& Induktion}
	
	Die Gesamtenergie des Magnetfelds einer stromdurchflossenen Spule beträgt:
	$E=\frac{1}{2}LI^{2}$
	
	Die magnetische Flussdichte B:
	$B=\mu\cdot\vec{H}$
	$\left[\SI{1}{\tesla}=\SI{1}{\newton\per\ampere\meter}=\SI{1}{\volt\second\per\meter\squared}\right]$
	$H=\frac{I}{2\pi r}$ $\left[\SI{1}{\ampere\per\meter}\right]$
	
	Die Stärke eines Feldes um einen stromdurchflossenen Leiter hängt hauptsächlich von der Stromstärke I ab.
	Es wird mit zunehmender Entfernung $r$ schwächer.
	$B=\mu_{0}\cdot\mu_{r}\cdot\frac{I}{2\pi r}$
	
	Die Stärke eines Feldes um eine Spule hängt neben der Stromstärke I auch von der Windungszahl N und der Länge $l$ ab.
	$B=\mu_{0}\cdot\mu_{r}\cdot\frac{N\cdot I}{l}$
	
	Induktion
	$U=+\frac{d\Phi}{dt}$
	wobei der magnetische Fluss
	$\Phi=\int_{A}\vec{B}\cdot d\vec{A}$
	ist.
	
	Vereinfachung für eine Leiterschleife:
	$U=-\frac{d}{dt}\int_{A}\vec{B}\cdot d\vec{A}$
	
	\section{Komplexe Wechselströme}
	
	\subsection{Effektivwerte}
	$I=\frac{\hat{i}}{\sqrt{2}}$
	$U=\frac{\hat{u}}{\sqrt{2}}$
	
	\subsection{Komplexe Darstellung von Wechselströmen}
	$u(t)=\hat{u}\sin(\omega t+\varphi_{u})=\hat{u}e^{j(\omega t+\varphi_{u})}$
	$i(t)=\hat{i}\sin(\omega t+\varphi_{i})=\hat{i}e^{j(\omega t+\varphi_{i})}$
	$\underline{Z_{C}}=\frac{1}{j\omega C}=-j\omega C$
	$\underline{Z_{L}}=j\omega L$
	
	\subsection{Komplexe Leistung}
	
	Momentanleistung
	$p(t)=U\cdot I\cdot \cos\varphi-U\cdot I\cdot \cos(2\omega t+\varphi_{u}+\varphi_{i})$
	$I(t)=\frac{\overline{U}}{\omega L}\sin(\omega\cdot t+\varphi-\frac{\pi}{2})$
	
	Mittlere Leistung:
	$P=U\cdot I\cdot \cos~\varphi$
	$\varphi=\varphi_{i}-\varphi_{u}$
	
	Scheinleistung:
	$S=U\cdot I$
	
	Blindleistung:
	$Q=U\cdot I~\sin\varphi=XI^{2}$
	
	Wirkleistung:
	$P=U\cdot I~\cos\varphi=RI^{2}$
	
	\subsection{Komplexer Widerstand}
	
	Komplexer Widerstand:
	$\underline{Z}=R+jX$
	
	Blindwiderstand:
	$X_{L}=\omega L$
	$X_{C}=-\frac{1}{\omega C}$
	
\end{document}